%**************************************************************************************
% บทคัดย่อภาษาไทย
%**************************************************************************************
\phantomsection
\addcontentsline{toc}{chapter}{บทคัดย่อภาษาไทย}
\vspace*{-0.5in}

\begin{tabularx}{\textwidth}{lX}
ชื่อ                           & : นายฐิติพงษ์ ตะเพียนทอง         %ชื่อผู้จัดทำ1 
%\\     %เว้นบรรทัด เมื่อมีเพื่อนร่วมโปรเจ็ค
%                             & : นางสาวสิวาลัย จินเจือ            %ชื่อผู้จัดทำ2 
\\
ชื่อปริญญานิพนธ์                  & : \hangindent=0.5em          % มีไว้สำหรับตอนชื่อโปรเจ็คเกิน 1 บรรทัด
                                แพลตฟอร์มคลัสเตอร์เพื่อสนับสนุนการเรียนการสอนและงานวิชาการภาควิชาเทคโนโลยีสารสนเทศ  %ระบุชื่อโปรเจ็คภาษาไทย
\\
% สาขาวิชา                       & : เทคโนโลยีสารสนเทศ % กรณีเป็นนักศึกษาหลักสูตร IT
%สาขาวิชา                       & : เทคโนโลยีสารสนเทศ (ต่อเนื่อง) % กรณีเป็นนักศึกษาหลักสูตร ITI
สาขาวิชา                       & : วิศวกรรมสารสนเทศและเครือข่าย % กรณีเป็นนักศึกษาหลักสูตร INE และ INET
\\
                             & : มหาวิทยาลัยเทคโนโลยีพระจอมเกล้าพระนครเหนือ
\\
อาจารย์ที่ปรึกษาปริญญานิพนธ์        & : ผู้ช่วยศาสตราจารย์ ดร.ขนิษฐา นามี    %ระบุชื่ออาจารย์ที่ปรึกษา
%\\      %เว้นบรรทัด เมื่อมีที่ปรึกษาร่วม
%ที่ปรึกษาร่วม                    & : ผู้ช่วยศาสตาจารย์ ดร.นิติการ นาคเจือทอง    %ระบุชื่อที่ปรึกษาร่วม (ถ้ามี)
\\
ปีการศึกษา                     & : 2568     %ระบุปีการศึกษา
\end{tabularx}

\vspace{5mm}
\phantomsection
\begin{center}\textbf{บทคัดย่อ}\end{center}
\hspace*{1.5em} %ย่อหน้า
ปริญญานิพนธ์ "แพลตฟอร์มคลัสเตอร์เพื่อสนับสนุนการเรียนการสอนและงานวิชาการภาควิชาเทคโนโลยีสารสนเทศ" ฉบับนี้ มีวัตถุประสงค์เพื่อพัฒนาแพลตฟอร์มคลัสเตอร์เสมือนที่สามารถสร้างและจัดการเครื่องเสมือนลีนุกซ์ได้อย่างรวดเร็วและมีประสิทธิภาพ เพื่อสนับสนุนการเรียนการสอนและการทำโครงงานในภาควิชาเทคโนโลยีสารสนเทศ และพัฒนาเว็บแอปพลิเคชันแบบ Role-Based Access Control (RBAC) ที่รองรับการใช้งานของทั้งนักศึกษาและอาจารย์ โดยคำนึงถึงความง่ายในการใช้งาน เสถียรภาพของระบบ และความปลอดภัยในการเข้าถึง ระบบประกอบด้วยสามส่วนหลัก ได้แก่ ส่วนติดต่อผู้ใช้งาน บริการ API สำหรับจัดการการยืนยันตัวตน การจัดการผู้ใช้ และเชื่อมต่อระหว่างส่วนติดต่อผู้ใช้และบริการจัดการเครื่องเสมือน และบริการจัดการเครื่องเสมือนที่เชื่อมต่อกับระบบ Hypervisor เพื่อสร้าง กำหนดค่า และจัดการเครื่องเสมือน ระบบสามารถลดความซับซ้อนในการตั้งค่าสภาพแวดล้อมสำหรับการใช้งานระบบลีนุกซ์ เพิ่มประสิทธิภาพการใช้ทรัพยากรเซิร์ฟเวอร์ และสนับสนุนการเรียนการสอนที่เกี่ยวข้องกับระบบเซิร์ฟเวอร์ การพัฒนาแอปพลิเคชัน และการจัดการโครงสร้างพื้นฐานด้านไอทีได้อย่างมีประสิทธิภาพ

\begin{flushright}
(ปริญญานิพนธ์มีจำนวนทั้งสิ้น \pageref{LastPage} หน้า)

\vfill
\rule{10 cm}{0.4pt} อาจารย์ที่ปรึกษาปริญญานิพนธ์ 
%\\ \vspace{3mm}                        %เว้นกรณีมีที่ปรึกษาร่วม
%\rule{12.6 cm}{0.4pt} ที่ปรึกษาร่วม         % กรณีมีที่ปรึกษาร่วม

\end{flushright}





%**************************************************************************************
% บทคัดย่อภาษาอังกฤษ
%**************************************************************************************
\newpage
\phantomsection
\addcontentsline{toc}{chapter}{บทคัดย่อภาษาอังกฤษ}
\vspace*{-0.5in}

\begin{tabularx}{\textwidth}{lX}
Name                    & : Mr.Thitipong Tapianthong  %ชื่อผู้จัดทำ1
%\\      %เว้นบรรทัดเมื่อมีเพื่อนร่วมโปรเจ็ค
%                       & : Ms.Siwalai Chinchua      %ชื่อผู้จัดทำ2
\\
Project Title           & : \hangindent=0.5em % มีไว้สำหรับตอนชื่อโปรเจ็คเกิน 1 บรรทัด
                            Cluster-Based Platform for Supporting Teaching and
                            Academic Activities in the Department of Information Technology  %ระบุชื่อโปรเจ็คภาษาอังกฤษ
\\
% Major Field             & : Information Technology % กรณีเป็นนักศึกษาหลักสูตร IT
%Major Field             & : Information Technology (Continuing Program) % กรณีเป็นนักศึกษาหลักสูตร ITI
Major Field             & : Information and Network Engineering % กรณีเป็นนักศึกษาหลักสูตร INE และ INET
\\
                        & : King Mongkut’s University of Technology North Bangkok
\\
Project Advisor         & : Asst. Prof. Dr.Khanista Namee   %ระบุชื่ออาจารย์ที่ปรึกษา
%\\      %เว้นบรรทัด เมื่อมีที่ปรึกษาร่วม
%Co-Advisor              & : Asst. Prof. Dr.Nitigan Nakjuatong   %ระบุชื่อที่ปรึกษาร่วม (ถ้ามี)
\\
Academic Year           & : 2025    %ระบุปีการศึกษา
\end{tabularx}

\vspace{5mm}
\begin{center}\textbf{Abstract}\end{center}

\hspace*{1.5em} %ย่อหน้า
This project, "Cluster-Based Platform for Supporting Teaching and Academic Activities in the Department of Information Technology," aims to develop a virtual cluster platform capable of rapidly and efficiently creating and managing Linux virtual machines to support teaching and project work in the Department of Information Technology. The project also focuses on developing a Role-Based Access Control (RBAC) web application that supports both students and faculty, emphasizing ease of use, system stability, and access security. The system consists of three main components: a user interface, an API service for handling authentication and user management and communication between the frontend and VM operations service, and a VM management service that interfaces with the Hypervisor system to create, configure, and manage virtual machines. The system reduces the complexity of setting up Linux environments, improves server resource utilization efficiency, and effectively supports teaching related to server systems, application development, and IT infrastructure management.

\begin{flushright}
(Total \pageref{LastPage} pages)

\vfill

\rule{12 cm}{0.4pt} Project Advisor
%\\ \vspace{3mm}                          %เว้นกรณีมีที่ปรึกษาร่วม
%\rule{12.7 cm}{0.4pt} Co-Advisor         % กรณีมีที่ปรึกษาร่วม

\end{flushright}


%**************************************************************************************
% กิตติกรรมประกาศ
%**************************************************************************************
\newpage
\phantomsection
\addcontentsline{toc}{chapter}{กิตติกรรมประกาศ}
\vspace*{-0.5in}

\begin{center}\textbf{กิตติกรรมประกาศ}\end{center}

\hspace*{1.5em}  %ย่อหน้า
ปริญญานิพนธ์ "แพลตฟอร์มคลัสเตอร์เพื่อสนับสนุนการเรียนการสอนและงานวิชาการภาควิชาเทคโนโลยีสารสนเทศ" จะไม่สามารถสำเร็จลุล่วงไปได้ด้วยดีหากไม่ได้รับการช่วยเหลือและความอนุเคราะห์จากบุคคลหลายท่าน ทางผู้จัดทำขอขอบพระคุณ ผู้ช่วยศาสตราจารย์ ดร.ขนิษฐา นามี ซึ่งเป็นอาจารย์ที่ปรึกษา ที่ได้กรุณาให้คำปรึกษาแนะนำ และให้แนวทางในการพัฒนาแพลตฟอร์มคลัสเตอร์และการจัดทำปริญญานิพนธ์ รวมทั้งคณาจารย์ภาควิชาเทคโนโลยีสารสนเทศที่ได้ให้ความรู้ด้านเทคโนโลยีและการพัฒนาระบบให้กับผู้จัดทำเพื่อนำมาประยุกต์ในการพัฒนาโครงงานนี้ อีกทั้งขอขอบคุณภาควิชาเทคโนโลยีสารสนเทศที่ให้การสนับสนุนโครงสร้างพื้นฐานและทรัพยากรที่จำเป็นสำหรับการพัฒนาและทดสอบระบบ 
\\
\hspace*{1.5em}  %ย่อหน้า
ท้ายที่สุดนี้ ทางผู้จัดทำขอน้อมรำลึกพระคุณบิดา มารดา ผู้ซึ่งมีพระคุณอย่างสูงสุดที่ให้ความอุปการะผู้จัดทำมาโดยตลอด รวมทั้งผู้มีพระคุณทุกท่าน คณะผู้จัดทำรู้สึกซาบซึ้งเป็นอย่างยิ่ง จึงใคร่ขอขอบพระคุณเป็นอย่างสูงมา ณ โอกาสนี้ด้วย
\\
\begin{flushright}
\makebox[6cm][l]{{ฐิติพงษ์ ตะเพียนทอง}}      %ชื่อผู้จัดทำ1
%\\     %เว้นบรรทัด เมื่อมีเพื่อนร่วมโปรเจ็ค
%\makebox[6cm][l]{สิวาลัย จินเจือ}         %ชื่อผู้จัดทำ2

\end{flushright}